\section{構造物の強度・安全率}

\subsection{平成23年度 問1-5}
\label{q:h23-reg-basic-1-5}

\noindent\textbf{原典：} sources/H23/basic.pdf、PDF物理4頁。

\textbf{公式正答：} 4（single）。

\par\textit{原典本文の校正と独自解説は未作成。分類のみの登録項目。}

\subsection{平成27年度 問1-3}
\label{q:h27-reg-basic-1-3}

\noindent\textbf{原典：} sources/H27/basic.pdf、PDF物理2頁。

\textbf{公式正答：} 5（single）。

\par\textit{原典本文の校正と独自解説は未作成。分類のみの登録項目。}

\subsection{平成29年度 問1-2}
\label{q:h29-reg-basic-1-2}

\noindent\textbf{原典：} sources/H29/basic.pdf、PDF物理1頁。

\textbf{公式正答：} 5（single）。

\par\textit{原典本文の校正と独自解説は未作成。分類のみの登録項目。}

\subsection{令和7年度 問1-3}
\label{q:r07-reg-basic-1-3}

\noindent\textbf{原典：} sources/R07/basic.pdf、PDF物理2頁。

\textbf{公式正答：} 2（single）。

\par\textit{原典本文の校正と独自解説は未作成。分類のみの登録項目。}

\section{信頼性と故障}

\subsection{平成23年度 問1-1}
\label{q:h23-reg-basic-1-1}

\noindent\textbf{原典：} sources/H23/basic.pdf、PDF物理1頁。

\textbf{公式正答：} 2（single）。

\par\textit{原典本文の校正と独自解説は未作成。分類のみの登録項目。}

\subsection{平成24年度 問1-1}
\label{q:h24-reg-basic-1-1}

\noindent\textbf{原典：} sources/H24/basic.pdf、PDF物理1頁。

\textbf{公式正答：} 4（single）。

\par\textit{原典本文の校正と独自解説は未作成。分類のみの登録項目。}

\subsection{平成26年度 問1-4}
\label{q:h26-reg-basic-1-4}

\noindent\textbf{原典：} sources/H26/basic.pdf、PDF物理3頁。

\textbf{公式正答：} 3（single）。

\par\textit{原典本文の校正と独自解説は未作成。分類のみの登録項目。}

\subsection{平成26年度 問1-5}
\label{q:h26-reg-basic-1-5}

\noindent\textbf{原典：} sources/H26/basic.pdf、PDF物理3頁。

\textbf{公式正答：} 1（single）。

\par\textit{原典本文の校正と独自解説は未作成。分類のみの登録項目。}

\subsection{平成27年度 問1-1}
\label{q:h27-reg-basic-1-1}

\noindent\textbf{原典：} sources/H27/basic.pdf、PDF物理1頁。

\textbf{公式正答：} 5（single）。

\par\textit{原典本文の校正と独自解説は未作成。分類のみの登録項目。}

\subsection{平成28年度 問1-1}
\label{q:h28-reg-basic-1-1}

\noindent\textbf{原典：} sources/H28/basic.pdf、PDF物理1頁。

\textbf{公式正答：} 4（single）。

\par\textit{原典本文の校正と独自解説は未作成。分類のみの登録項目。}

\subsection{平成30年度 問1-1}
\label{q:h30-reg-basic-1-1}

\noindent\textbf{原典：} sources/H30/basic.pdf、PDF物理1頁。

\textbf{公式正答：} 3（single）。

\par\textit{原典本文の校正と独自解説は未作成。分類のみの登録項目。}

\subsection{令和元年度再試験 問1-3}
\label{q:r01-ret-basic-1-3}

\noindent\textbf{原典：} sources/R01\_retest/basic.pdf、PDF物理3頁。

\textbf{公式正答：} 1（single）。

\par\textit{原典本文の校正と独自解説は未作成。分類のみの登録項目。}

\subsection{令和元年度再試験 問1-6}
\label{q:r01-ret-basic-1-6}

\noindent\textbf{原典：} sources/R01\_retest/basic.pdf、PDF物理5頁。

\textbf{公式正答：} 5（single）。

\par\textit{原典本文の校正と独自解説は未作成。分類のみの登録項目。}

\subsection{令和2年度 問1-6}
\label{q:r02-reg-basic-1-6}

\noindent\textbf{原典：} sources/R02/basic.pdf、PDF物理6頁。

\textbf{公式正答：} 3（single）。

\par\textit{原典本文の校正と独自解説は未作成。分類のみの登録項目。}

\subsection{令和3年度 問1-2}
\label{q:r03-reg-basic-1-2}

\noindent\textbf{原典：} sources/R03/basic.pdf、PDF物理2頁。

\textbf{公式正答：} 2（single）。

\par\textit{原典本文の校正と独自解説は未作成。分類のみの登録項目。}

\subsection{令和3年度 問1-4}
\label{q:r03-reg-basic-1-4}

\noindent\textbf{原典：} sources/R03/basic.pdf、PDF物理3頁。

\textbf{公式正答：} 3（single）。

\par\textit{原典本文の校正と独自解説は未作成。分類のみの登録項目。}

\subsection{令和5年度 問1-4}
\label{q:r05-reg-basic-1-4}

\noindent\textbf{原典：} sources/R05/basic.pdf、PDF物理3頁。

\textbf{公式正答：} 4（single）。

\par\textit{原典本文の校正と独自解説は未作成。分類のみの登録項目。}

\subsection{令和5年度 問1-5}
\label{q:r05-reg-basic-1-5}

\noindent\textbf{原典：} sources/R05/basic.pdf、PDF物理4頁。

\textbf{公式正答：} 1（single）。

\par\textit{原典本文の校正と独自解説は未作成。分類のみの登録項目。}

\section{待ち行列}

\subsection{平成23年度 問1-2}
\label{q:h23-reg-basic-1-2}

\noindent\textbf{原典：} sources/H23/basic.pdf、PDF物理1頁。

\textbf{公式正答：} 4（single）。

\par\textit{原典本文の校正と独自解説は未作成。分類のみの登録項目。}

\subsection{平成25年度 問1-5}
\label{q:h25-reg-basic-1-5}

\noindent\textbf{原典：} sources/H25/basic.pdf、PDF物理5頁。

\textbf{公式正答：} 4（single）。

\par\textit{原典本文の校正と独自解説は未作成。分類のみの登録項目。}

\subsection{平成27年度 問1-2}
\label{q:h27-reg-basic-1-2}

\noindent\textbf{原典：} sources/H27/basic.pdf、PDF物理2頁。

\textbf{公式正答：} 4（single）。

\par\textit{原典本文の校正と独自解説は未作成。分類のみの登録項目。}

\subsection{平成29年度 問1-1}
\label{q:h29-reg-basic-1-1}

\noindent\textbf{原典：} sources/H29/basic.pdf、PDF物理1頁。

\textbf{公式正答：} 5（single）。

\par\textit{原典本文の校正と独自解説は未作成。分類のみの登録項目。}

\subsection{令和元年度 問1-5}
\label{q:r01-reg-basic-1-5}

\noindent\textbf{原典：} sources/R01/basic.pdf、PDF物理5頁。

\textbf{公式正答：} 2（single）。

\par\textit{原典本文の校正と独自解説は未作成。分類のみの登録項目。}

\subsection{令和7年度 問1-5}
\label{q:r07-reg-basic-1-5}

\noindent\textbf{原典：} sources/R07/basic.pdf、PDF物理3頁。

\textbf{公式正答：} 5（single）。

\par\textit{原典本文の校正と独自解説は未作成。分類のみの登録項目。}

\section{工程管理・PERT}

\subsection{平成23年度 問1-3}
\label{q:h23-reg-basic-1-3}

\noindent\textbf{原典：} sources/H23/basic.pdf、PDF物理2頁。

\textbf{公式正答：} 3（single）。

\par\textit{原典本文の校正と独自解説は未作成。分類のみの登録項目。}

\subsection{平成25年度 問1-4}
\label{q:h25-reg-basic-1-4}

\noindent\textbf{原典：} sources/H25/basic.pdf、PDF物理4頁。

\textbf{公式正答：} 4（single）。

\par\textit{原典本文の校正と独自解説は未作成。分類のみの登録項目。}

\subsection{平成28年度 問1-4}
\label{q:h28-reg-basic-1-4}

\noindent\textbf{原典：} sources/H28/basic.pdf、PDF物理4頁。

\textbf{公式正答：} 4（single）。

\par\textit{原典本文の校正と独自解説は未作成。分類のみの登録項目。}

\subsection{平成30年度 問1-2}
\label{q:h30-reg-basic-1-2}

\noindent\textbf{原典：} sources/H30/basic.pdf、PDF物理2頁。

\textbf{公式正答：} 1（single）。

\par\textit{原典本文の校正と独自解説は未作成。分類のみの登録項目。}

\subsection{令和元年度再試験 問1-4}
\label{q:r01-ret-basic-1-4}

\noindent\textbf{原典：} sources/R01\_retest/basic.pdf、PDF物理4頁。

\textbf{公式正答：} 2（single）。

\par\textit{原典本文の校正と独自解説は未作成。分類のみの登録項目。}

\subsection{令和3年度 問1-3}
\label{q:r03-reg-basic-1-3}

\noindent\textbf{原典：} sources/R03/basic.pdf、PDF物理3頁。

\textbf{公式正答：} 2（single）。

\par\textit{原典本文の校正と独自解説は未作成。分類のみの登録項目。}

\subsection{令和7年度 問1-2}
\label{q:r07-reg-basic-1-2}

\noindent\textbf{原典：} sources/R07/basic.pdf、PDF物理2頁。

\textbf{公式正答：} 4（single）。

\par\textit{原典本文の校正と独自解説は未作成。分類のみの登録項目。}

\section{最適化・意思決定}

\subsection{平成23年度 問1-4}
\label{q:h23-reg-basic-1-4}

\noindent\textbf{原典：} sources/H23/basic.pdf、PDF物理3頁。

\textbf{公式正答：} 1（single）。

\par\textit{原典本文の校正と独自解説は未作成。分類のみの登録項目。}

\subsection{平成24年度 問1-5}
\label{q:h24-reg-basic-1-5}

\noindent\textbf{原典：} sources/H24/basic.pdf、PDF物理4頁。

\textbf{公式正答：} 5（single）。

\par\textit{原典本文の校正と独自解説は未作成。分類のみの登録項目。}

\subsection{平成25年度 問1-3}
\label{q:h25-reg-basic-1-3}

\noindent\textbf{原典：} sources/H25/basic.pdf、PDF物理3頁。

\textbf{公式正答：} 2（single）。

\par\textit{原典本文の校正と独自解説は未作成。分類のみの登録項目。}

\subsection{平成26年度 問1-2}
\label{q:h26-reg-basic-1-2}

\noindent\textbf{原典：} sources/H26/basic.pdf、PDF物理1頁。

\textbf{公式正答：} 4（single）。

\par\textit{原典本文の校正と独自解説は未作成。分類のみの登録項目。}

\subsection{平成28年度 問1-5}
\label{q:h28-reg-basic-1-5}

\noindent\textbf{原典：} sources/H28/basic.pdf、PDF物理5頁。

\textbf{公式正答：} 5（single）。

\par\textit{原典本文の校正と独自解説は未作成。分類のみの登録項目。}

\subsection{平成29年度 問1-3}
\label{q:h29-reg-basic-1-3}

\noindent\textbf{原典：} sources/H29/basic.pdf、PDF物理2頁。

\textbf{公式正答：} 2（single）。

\par\textit{原典本文の校正と独自解説は未作成。分類のみの登録項目。}

\subsection{平成30年度 問1-4}
\label{q:h30-reg-basic-1-4}

\noindent\textbf{原典：} sources/H30/basic.pdf、PDF物理4頁。

\textbf{公式正答：} 2（single）。

\par\textit{原典本文の校正と独自解説は未作成。分類のみの登録項目。}

\subsection{令和元年度 問1-1}
\label{q:r01-reg-basic-1-1}

\noindent\textbf{原典：} sources/R01/basic.pdf、PDF物理1頁。

\textbf{公式正答：} 5（single）。

\par\textit{原典本文の校正と独自解説は未作成。分類のみの登録項目。}

\subsection{令和元年度再試験 問1-2}
\label{q:r01-ret-basic-1-2}

\noindent\textbf{原典：} sources/R01\_retest/basic.pdf、PDF物理2頁。

\textbf{公式正答：} 3（single）。

\par\textit{原典本文の校正と独自解説は未作成。分類のみの登録項目。}

\subsection{令和2年度 問1-4}
\label{q:r02-reg-basic-1-4}

\noindent\textbf{原典：} sources/R02/basic.pdf、PDF物理4頁。

\textbf{公式正答：} 5（single）。

\par\textit{原典本文の校正と独自解説は未作成。分類のみの登録項目。}

\subsection{令和6年度 問1-5}
\label{q:r06-reg-basic-1-5}

\noindent\textbf{原典：} sources/R06/basic.pdf、PDF物理4頁。

\textbf{公式正答：} 4（single）。

\par\textit{原典本文の校正と独自解説は未作成。分類のみの登録項目。}

\section{品質管理・統計}

\subsection{平成24年度 問1-3}
\label{q:h24-reg-basic-1-3}

\noindent\textbf{原典：} sources/H24/basic.pdf、PDF物理3頁。

\textbf{公式正答：} 3（single）。

\par\textit{原典本文の校正と独自解説は未作成。分類のみの登録項目。}

\subsection{平成25年度 問1-6}
\label{q:h25-reg-basic-1-6}

\noindent\textbf{原典：} sources/H25/basic.pdf、PDF物理5頁。

\textbf{公式正答：} 3（single）。

\par\textit{原典本文の校正と独自解説は未作成。分類のみの登録項目。}

\subsection{平成26年度 問1-3}
\label{q:h26-reg-basic-1-3}

\noindent\textbf{原典：} sources/H26/basic.pdf、PDF物理2頁。

\textbf{公式正答：} 1（single）。

\par\textit{原典本文の校正と独自解説は未作成。分類のみの登録項目。}

\subsection{平成27年度 問1-5}
\label{q:h27-reg-basic-1-5}

\noindent\textbf{原典：} sources/H27/basic.pdf、PDF物理4頁。

\textbf{公式正答：} 1（single）。

\par\textit{原典本文の校正と独自解説は未作成。分類のみの登録項目。}

\subsection{平成27年度 問1-6}
\label{q:h27-reg-basic-1-6}

\noindent\textbf{原典：} sources/H27/basic.pdf、PDF物理4頁。

\textbf{公式正答：} 5（single）。

\par\textit{原典本文の校正と独自解説は未作成。分類のみの登録項目。}

\subsection{平成28年度 問1-2}
\label{q:h28-reg-basic-1-2}

\noindent\textbf{原典：} sources/H28/basic.pdf、PDF物理2頁。

\textbf{公式正答：} 4（single）。

\par\textit{原典本文の校正と独自解説は未作成。分類のみの登録項目。}

\subsection{平成30年度 問1-5}
\label{q:h30-reg-basic-1-5}

\noindent\textbf{原典：} sources/H30/basic.pdf、PDF物理5頁。

\textbf{公式正答：} 3（single）。

\par\textit{原典本文の校正と独自解説は未作成。分類のみの登録項目。}

\subsection{令和元年度 問1-3}
\label{q:r01-reg-basic-1-3}

\noindent\textbf{原典：} sources/R01/basic.pdf、PDF物理3頁。

\textbf{公式正答：} 5（single）。

\par\textit{原典本文の校正と独自解説は未作成。分類のみの登録項目。}

\subsection{令和2年度 問1-2}
\label{q:r02-reg-basic-1-2}

\noindent\textbf{原典：} sources/R02/basic.pdf、PDF物理2頁。

\textbf{公式正答：} 4（single）。

\par\textit{原典本文の校正と独自解説は未作成。分類のみの登録項目。}

\subsection{令和4年度 問1-3}
\label{q:r04-reg-basic-1-3}

\noindent\textbf{原典：} sources/R04/basic.pdf、PDF物理2頁。

\textbf{公式正答：} 3（single）。

\par\textit{原典本文の校正と独自解説は未作成。分類のみの登録項目。}

\subsection{令和5年度 問1-6}
\label{q:r05-reg-basic-1-6}

\noindent\textbf{原典：} sources/R05/basic.pdf、PDF物理4頁。

\textbf{公式正答：} 2（single）。

\par\textit{原典本文の校正と独自解説は未作成。分類のみの登録項目。}

\subsection{令和7年度 問1-6}
\label{q:r07-reg-basic-1-6}

\noindent\textbf{原典：} sources/R07/basic.pdf、PDF物理4頁。

\textbf{公式正答：} 3（single）。

\par\textit{原典本文の校正と独自解説は未作成。分類のみの登録項目。}

\section{論理式・命題}

\subsection{平成28年度 問1-6}
\label{q:h28-reg-basic-1-6}

\noindent\textbf{原典：} sources/H28/basic.pdf、PDF物理6頁。

\textbf{公式正答：} 1（single）。

\par\textit{原典本文の校正と独自解説は未作成。分類のみの登録項目。}

\section{数値表現}

\subsection{平成26年度 問1-6}
\label{q:h26-reg-basic-1-6}

\noindent\textbf{原典：} sources/H26/basic.pdf、PDF物理4頁。

\textbf{公式正答：} 3（single）。

\par\textit{原典本文の校正と独自解説は未作成。分類のみの登録項目。}

\subsection{平成29年度 問1-5}
\label{q:h29-reg-basic-1-5}

\noindent\textbf{原典：} sources/H29/basic.pdf、PDF物理4頁。

\textbf{公式正答：} 1（single）。

\par\textit{原典本文の校正と独自解説は未作成。分類のみの登録項目。}

\section{微積分・求積}

\subsection{令和元年度 問1-6}
\label{q:r01-reg-basic-1-6}

\noindent\textbf{原典：} sources/R01/basic.pdf、PDF物理6頁。

\textbf{公式正答：} 3（single）。

\par\textit{原典本文の校正と独自解説は未作成。分類のみの登録項目。}

\subsection{令和元年度再試験 問1-1}
\label{q:r01-ret-basic-1-1}

\noindent\textbf{原典：} sources/R01\_retest/basic.pdf、PDF物理1頁。

\textbf{公式正答：} 5（single）。

\par\textit{原典本文の校正と独自解説は未作成。分類のみの登録項目。}

\section{材料の応力・変形}

\subsection{令和2年度 問1-5}
\label{q:r02-reg-basic-1-5}

\noindent\textbf{原典：} sources/R02/basic.pdf、PDF物理5頁。

\textbf{公式正答：} 5（single）。

\par\textit{原典本文の校正と独自解説は未作成。分類のみの登録項目。}

\subsection{令和7年度 問1-4}
\label{q:r07-reg-basic-1-4}

\noindent\textbf{原典：} sources/R07/basic.pdf、PDF物理3頁。

\textbf{公式正答：} 1（single）。

\par\textit{原典本文の校正と独自解説は未作成。分類のみの登録項目。}

\section{環境法・循環・持続可能性}

\subsection{令和7年度 問1-1}
\label{q:r07-reg-basic-1-1}

\noindent\textbf{原典：} sources/R07/basic.pdf、PDF物理1頁。

\textbf{公式正答：} 3（single）。

\par\textit{原典本文の校正と独自解説は未作成。分類のみの登録項目。}

\section{断面二次モーメント・断面係数と寸法依存}

\subsection{令和4年度 問1-5}
\label{q:r04-reg-basic-1-5}

\noindent\textbf{原典：} sources/R04/basic.pdf、PDF物理3頁。

\textbf{公式正答：} 3（single）。

\par\textbf{検討未確定：} 径/高さを変える比較でq,L,E（矩形は幅b）を固定する前提が必要。自重のみでqが断面積に比例して変わる円梁なら応力はdの逆1乗、たわみは逆2乗となる。試験の通常の固定荷重解と区別。

\paragraph{問題の要旨} 等分布荷重を受ける単純支持梁の断面寸法と、最大曲げ応力・たわみのべき依存性を導く。

\Needspace{4\baselineskip}\paragraph{確認した条件・記述（独自要約）}\mbox{}\par

\noindent\textbf{判断方向：} 空欄の適切な語・数値又は正誤の組合せ。全候補の順序は各肢に記録。\par

\noindent\textbf{原典照合：} 16原典画像で17問の全候補、図、数値、不等号、表の列順を確認。\par

\noindent\textbf{前提：} 解法に記録した弾性範囲、荷重条件、幾何、独立性を明示する。\par

\paragraph{問われる核心} 等分布荷重を受ける単純支持梁の断面寸法と、最大曲げ応力・たわみのべき依存性を導く。

\paragraph{必要知識} 固定するq,L,E、円断面IとZ、矩形の幅固定、曲げモーメント、梁の線形小たわみ。

\paragraph{考え方・正答を選ぶ理由（査読待ち初稿）} 同じq,L,Eの円断面梁では \(M_{\max}=qL^2/8\)、\(I=\pi d^4/64\)、\(Z=I/(d/2)=\pi d^3/32\)。従って \(\sigma_{\max}=M_{\max}/Z\propto d^{-3}\)、\(\delta_{\max}=5qL^4/(384EI)\propto d^{-4}\)。矩形断面は幅bを固定して \(Z=bh^2/6\) より \(\sigma_{\max}\propto h^{-2}\)。A3、B反比例、C4、D2の3となる。

\Needspace{4\baselineskip}\paragraph{選択肢の検討（査読待ち初稿）}\mbox{}\par

\noindent\textbf{1：不適合} 3・比例・4・3。列2は反比例であり比例ではない。列4は2であり3ではない。\par

\noindent\textbf{2：不適合} 4・比例・3・2。列1は3であり4ではない。列2は反比例であり比例ではない。列3は4であり3ではない。\par

\noindent\textbf{3：適合（前提内）} 3・反比例・4・2。曲げ応力は断面係数Zの逆数、変位は断面二次モーメントIの逆数という二つの依存性に適合する。\par

\noindent\textbf{4：不適合} 4・反比例・3・3。列1は3であり4ではない。列3は4であり3ではない。列4は2であり3ではない。\par

\noindent\textbf{5：不適合} 3・反比例・4・3。列4は2であり3ではない。\par

\paragraph{再利用する考え方} 材料の強さ、構造の安定、変形による機能低下を分ける。寸法だけを変えるとき何を固定したかを明示し、同じ理論でも候補番号を年度間で流用しない。

\paragraph{共通理論} \hyperref[theory:structural:beam]{第\ref*{theory:structural:beam}節}, \hyperref[theory:structural:scaling]{第\ref*{theory:structural:scaling}節}

\par 参考：\cite{struct_mit_beam,struct_mit_buckling,struct_mlit_limit}。

\paragraph{関連問題} \hyperref[q:r06-reg-basic-1-3]{R06-REG-BASIC-1-3}

\par 難易度 3（校準初稿）、優先度未評価。

\subsection{令和6年度 問1-3}
\label{q:r06-reg-basic-1-3}

\noindent\textbf{原典：} sources/R06/basic.pdf、PDF物理2頁。

\textbf{公式正答：} 1（single）。

\par\textbf{検討未確定：} 径/高さを変える比較でq,L,E（矩形は幅b）を固定する前提が必要。自重のみでqが断面積に比例して変わる円梁なら応力はdの逆1乗、たわみは逆2乗となる。試験の通常の固定荷重解と区別。

\paragraph{問題の要旨} 等分布荷重を受ける単純支持梁の断面寸法と、最大曲げ応力・たわみのべき依存性を導く。

\Needspace{4\baselineskip}\paragraph{確認した条件・記述（独自要約）}\mbox{}\par

\noindent\textbf{判断方向：} 空欄の適切な語・数値又は正誤の組合せ。全候補の順序は各肢に記録。\par

\noindent\textbf{原典照合：} 16原典画像で17問の全候補、図、数値、不等号、表の列順を確認。\par

\noindent\textbf{前提：} 解法に記録した弾性範囲、荷重条件、幾何、独立性を明示する。\par

\paragraph{問われる核心} 等分布荷重を受ける単純支持梁の断面寸法と、最大曲げ応力・たわみのべき依存性を導く。

\paragraph{必要知識} 固定するq,L,E、円断面IとZ、矩形の幅固定、曲げモーメント、梁の線形小たわみ。

\paragraph{考え方・正答を選ぶ理由（査読待ち初稿）} 同じq,L,Eの円断面単純支持梁で \(I=\pi d^4/64\)、\(Z=\pi d^3/32\) である。最大モーメントは \(qL^2/8\)、最大たわみは \(5qL^4/(384EI)\)。従って \(\sigma_{\max}\propto d^{-3}\)、\(\delta_{\max}\propto d^{-4}\) であり、3乗に反比例・4乗に反比例の1を選ぶ。

\Needspace{4\baselineskip}\paragraph{選択肢の検討（査読待ち初稿）}\mbox{}\par

\noindent\textbf{1：適合（前提内）} 3・反比例・4・反比例。曲げ応力は断面係数Zの逆数、変位は断面二次モーメントIの逆数という二つの依存性に適合する。\par

\noindent\textbf{2：不適合} 4・比例・4・反比例。列1は3であり4ではない。列2は反比例であり比例ではない。\par

\noindent\textbf{3：不適合} 4・反比例・3・反比例。列1は3であり4ではない。列3は4であり3ではない。\par

\noindent\textbf{4：不適合} 4・比例・3・比例。列1は3であり4ではない。列2は反比例であり比例ではない。列3は4であり3ではない。列4は反比例であり比例ではない。\par

\noindent\textbf{5：不適合} 3・反比例・4・比例。列4は反比例であり比例ではない。\par

\paragraph{再利用する考え方} 材料の強さ、構造の安定、変形による機能低下を分ける。寸法だけを変えるとき何を固定したかを明示し、同じ理論でも候補番号を年度間で流用しない。

\paragraph{共通理論} \hyperref[theory:structural:beam]{第\ref*{theory:structural:beam}節}, \hyperref[theory:structural:scaling]{第\ref*{theory:structural:scaling}節}

\par 参考：\cite{struct_mit_beam,struct_mit_buckling,struct_mlit_limit}。

\paragraph{関連問題} \hyperref[q:r04-reg-basic-1-5]{R04-REG-BASIC-1-5}

\par 難易度 3（校準初稿）、優先度未評価。

\section{圧縮安定とオイラー座屈}

\subsection{平成25年度 問1-2}
\label{q:h25-reg-basic-1-2}

\noindent\textbf{原典：} sources/H25/basic.pdf、PDF物理2頁。

\textbf{公式正答：} 5（single）。

\paragraph{問題の要旨} 薄い定規の引張と圧縮の差を座屈で説明し、許容応力と安全率を対応させる。

\Needspace{4\baselineskip}\paragraph{確認した条件・記述（独自要約）}\mbox{}\par

\noindent\textbf{判断方向：} 空欄の適切な語・数値又は正誤の組合せ。全候補の順序は各肢に記録。\par

\noindent\textbf{原典照合：} 16原典画像で17問の全候補、図、数値、不等号、表の列順を確認。\par

\noindent\textbf{前提：} 解法に記録した弾性範囲、荷重条件、幾何、独立性を明示する。\par

\paragraph{問われる核心} 薄い定規の引張と圧縮の差を座屈で説明し、許容応力と安全率を対応させる。

\paragraph{必要知識} 材料強度と座屈荷重の相違、圧縮の安定性、許容応力＝強度/安全率。

\paragraph{考え方・正答を選ぶ理由（査読待ち初稿）} 引張では直線形状が安定しやすいが、圧縮では小さな横たわみが増幅され座屈する。Aは引張、Bは圧縮、Cは座屈、Dは安全率である。強度を \(\sigma_{\mathrm{str}}\)、安全率を \(n>1\) とすると \(\sigma_{\mathrm{allow}}=\sigma_{\mathrm{str}}/n\)。圧縮の座屈荷重はこの材料強度だけで決まらない。表の並びに従い本問の候補番号を選ぶ。

\Needspace{4\baselineskip}\paragraph{選択肢の検討（査読待ち初稿）}\mbox{}\par

\noindent\textbf{1：不適合} 引張・圧縮・剥離・安全率。横へたわむ圧縮不安定は座屈であり、層がはがれる剥離ではない。\par

\noindent\textbf{2：不適合} 圧縮・引張・剥離・安全率。薄い定規は引張より圧縮の安定性を失いやすく、AとBが逆である。横へたわむ圧縮不安定は座屈であり、層がはがれる剥離ではない。\par

\noindent\textbf{3：不適合} 圧縮・引張・剥離・ヤング率。薄い定規は引張より圧縮の安定性を失いやすく、AとBが逆である。横へたわむ圧縮不安定は座屈であり、層がはがれる剥離ではない。強度を割る無次元比は安全率であり、応力単位のヤング率ではない。\par

\noindent\textbf{4：不適合} 圧縮・引張・座屈・ヤング率。薄い定規は引張より圧縮の安定性を失いやすく、AとBが逆である。強度を割る無次元比は安全率であり、応力単位のヤング率ではない。\par

\noindent\textbf{5：適合（前提内）} 引張・圧縮・座屈・安全率。引張・圧縮・座屈・安全率が対応する。形の安定と材料強度は別の照査である。\par

\paragraph{再利用する考え方} 材料の強さ、構造の安定、変形による機能低下を分ける。寸法だけを変えるとき何を固定したかを明示し、同じ理論でも候補番号を年度間で流用しない。

\paragraph{共通理論} \hyperref[theory:structural:buckling]{第\ref*{theory:structural:buckling}節}, \hyperref[theory:structural:allowable]{第\ref*{theory:structural:allowable}節}

\par 参考：\cite{struct_mit_beam,struct_mit_buckling,struct_mlit_limit}。

\paragraph{関連問題} \hyperref[q:h28-reg-basic-1-3]{H28-REG-BASIC-1-3}, \hyperref[q:r03-reg-basic-1-5]{R03-REG-BASIC-1-5}

\par 難易度 2（校準初稿）、優先度未評価。

\subsection{平成28年度 問1-3}
\label{q:h28-reg-basic-1-3}

\noindent\textbf{原典：} sources/H28/basic.pdf、PDF物理3頁。

\textbf{公式正答：} 4（single）。

\paragraph{問題の要旨} 薄い定規の引張と圧縮の差を座屈で説明し、許容応力と安全率を対応させる。

\Needspace{4\baselineskip}\paragraph{確認した条件・記述（独自要約）}\mbox{}\par

\noindent\textbf{判断方向：} 空欄の適切な語・数値又は正誤の組合せ。全候補の順序は各肢に記録。\par

\noindent\textbf{原典照合：} 16原典画像で17問の全候補、図、数値、不等号、表の列順を確認。\par

\noindent\textbf{前提：} 解法に記録した弾性範囲、荷重条件、幾何、独立性を明示する。\par

\paragraph{問われる核心} 薄い定規の引張と圧縮の差を座屈で説明し、許容応力と安全率を対応させる。

\paragraph{必要知識} 材料強度と座屈荷重の相違、圧縮の安定性、許容応力＝強度/安全率。

\paragraph{考え方・正答を選ぶ理由（査読待ち初稿）} 引張では直線形状が安定しやすいが、圧縮では小さな横たわみが増幅され座屈する。Aは引張、Bは圧縮、Cは座屈、Dは安全率である。強度を \(\sigma_{\mathrm{str}}\)、安全率を \(n>1\) とすると \(\sigma_{\mathrm{allow}}=\sigma_{\mathrm{str}}/n\)。圧縮の座屈荷重はこの材料強度だけで決まらない。表の並びに従い本問の候補番号を選ぶ。

\Needspace{4\baselineskip}\paragraph{選択肢の検討（査読待ち初稿）}\mbox{}\par

\noindent\textbf{1：不適合} 引張・圧縮・剥離・安全率。横へたわむ圧縮不安定は座屈であり、層がはがれる剥離ではない。\par

\noindent\textbf{2：不適合} 圧縮・引張・剥離・安全率。薄い定規は引張より圧縮の安定性を失いやすく、AとBが逆である。横へたわむ圧縮不安定は座屈であり、層がはがれる剥離ではない。\par

\noindent\textbf{3：不適合} 圧縮・引張・剥離・ヤング率。薄い定規は引張より圧縮の安定性を失いやすく、AとBが逆である。横へたわむ圧縮不安定は座屈であり、層がはがれる剥離ではない。強度を割る無次元比は安全率であり、応力単位のヤング率ではない。\par

\noindent\textbf{4：適合（前提内）} 引張・圧縮・座屈・安全率。引張・圧縮・座屈・安全率が対応する。形の安定と材料強度は別の照査である。\par

\noindent\textbf{5：不適合} 圧縮・引張・座屈・ヤング率。薄い定規は引張より圧縮の安定性を失いやすく、AとBが逆である。強度を割る無次元比は安全率であり、応力単位のヤング率ではない。\par

\paragraph{再利用する考え方} 材料の強さ、構造の安定、変形による機能低下を分ける。寸法だけを変えるとき何を固定したかを明示し、同じ理論でも候補番号を年度間で流用しない。

\paragraph{共通理論} \hyperref[theory:structural:buckling]{第\ref*{theory:structural:buckling}節}, \hyperref[theory:structural:allowable]{第\ref*{theory:structural:allowable}節}

\par 参考：\cite{struct_mit_beam,struct_mit_buckling,struct_mlit_limit}。

\paragraph{関連問題} \hyperref[q:h25-reg-basic-1-2]{H25-REG-BASIC-1-2}

\par 難易度 2（校準初稿）、優先度未評価。

\subsection{令和元年度 問1-4}
\label{q:r01-reg-basic-1-4}

\noindent\textbf{原典：} sources/R01/basic.pdf、PDF物理4頁。

\textbf{公式正答：} 5（single）。

\paragraph{問題の要旨} 細長い棒の軸圧縮による横変形を、静的座屈として説明する。

\Needspace{4\baselineskip}\paragraph{確認した条件・記述（独自要約）}\mbox{}\par

\noindent\textbf{判断方向：} 空欄の適切な語・数値又は正誤の組合せ。全候補の順序は各肢に記録。\par

\noindent\textbf{原典照合：} 16原典画像で17問の全候補、図、数値、不等号、表の列順を確認。\par

\noindent\textbf{前提：} 解法に記録した弾性範囲、荷重条件、幾何、独立性を明示する。\par

\paragraph{問われる核心} 細長い棒の軸圧縮による横変形を、静的座屈として説明する。

\paragraph{必要知識} 圧縮荷重、曲げ剛性EI、端部拘束に依存する有効長KL、座屈と共振の差。

\paragraph{考え方・正答を選ぶ理由（査読待ち初稿）} 軸圧縮のもとで横たわみを増やす効果が曲げの復元に勝つと直線状態が不安定になる。\(P_{\mathrm{cr}}=\pi^2EI/(KL)^2\) であり、断面円の \(I=\pi d^4/64\) より、小径ほど臨界荷重が小さい。図の現象はねじれや振動の共振ではなく曲げ座屈である。理想棒では臨界を超える条件を選び、実物の初期不整に対しては臨界前からたわみが生じ得る点も分ける。

\Needspace{4\baselineskip}\paragraph{選択肢の検討（査読待ち初稿）}\mbox{}\par

\noindent\textbf{1：不適合} 小さい・以下・ねじれ・座屈。理想的な直線平衡の安定性を失うのは臨界荷重を超える側である。描かれた横変位は曲げであり、断面の軸回り回転であるねじれではない。\par

\noindent\textbf{2：不適合} 大きい・以下・ねじれ・共振。同じ長さ・材質では小径ほど曲げ剛性EIが小さく座屈しやすい。理想的な直線平衡の安定性を失うのは臨界荷重を超える側である。描かれた横変位は曲げであり、断面の軸回り回転であるねじれではない。静的圧縮での安定性喪失は座屈であり、周期的外力と固有振動の共振ではない。\par

\noindent\textbf{3：不適合} 小さい・超える・ねじれ・共振。描かれた横変位は曲げであり、断面の軸回り回転であるねじれではない。静的圧縮での安定性喪失は座屈であり、周期的外力と固有振動の共振ではない。\par

\noindent\textbf{4：不適合} 大きい・超える・曲げ・共振。同じ長さ・材質では小径ほど曲げ剛性EIが小さく座屈しやすい。静的圧縮での安定性喪失は座屈であり、周期的外力と固有振動の共振ではない。\par

\noindent\textbf{5：適合（前提内）} 小さい・超える・曲げ・座屈。小さい径・臨界荷重を超える・曲げ・座屈の対応である。\par

\paragraph{再利用する考え方} 材料の強さ、構造の安定、変形による機能低下を分ける。寸法だけを変えるとき何を固定したかを明示し、同じ理論でも候補番号を年度間で流用しない。

\paragraph{共通理論} \hyperref[theory:structural:buckling]{第\ref*{theory:structural:buckling}節}

\par 参考：\cite{struct_mit_beam,struct_mit_buckling,struct_mlit_limit}。

\paragraph{関連問題} \hyperref[q:r05-reg-basic-1-2]{R05-REG-BASIC-1-2}

\par 難易度 2（校準初稿）、優先度未評価。

\subsection{令和3年度 問1-5}
\label{q:r03-reg-basic-1-5}

\noindent\textbf{原典：} sources/R03/basic.pdf、PDF物理4頁。

\textbf{公式正答：} 5（single）。

\paragraph{問題の要旨} オイラー座屈の長さ・強度依存性と、許容応力・限界状態を組合せで検討する。

\Needspace{4\baselineskip}\paragraph{確認した条件・記述（独自要約）}\mbox{}\par

\noindent\textbf{判断方向：} 空欄の適切な語・数値又は正誤の組合せ。全候補の順序は各肢に記録。\par

\noindent\textbf{原典照合：} 16原典画像で17問の全候補、図、数値、不等号、表の列順を確認。\par

\noindent\textbf{前提：} 解法に記録した弾性範囲、荷重条件、幾何、独立性を明示する。\par

\paragraph{問われる核心} オイラー座屈の長さ・強度依存性と、許容応力・限界状態を組合せで検討する。

\paragraph{必要知識} オイラー式の弾性条件、Lの逆二乗、ヤング率と降伏強度の差、安全率の除算、複数の限界状態。

\paragraph{考え方・正答を選ぶ理由（査読待ち初稿）} \(P_{\mathrm{cr}}\propto L^{-2}\) なので長い柱ほど座屈荷重は小さい。弾性座屈域なら同じ \(E,I,KL\) の柱の臨界荷重は降伏強度を直接含まない。ただし算出座屈応力が降伏域に入る柱には弾性式をそのまま適用しない。許容応力は強度を安全率で割る。使用・終局・疲労等の限界状態を分けるため、誤正誤誤の5である。

\Needspace{4\baselineskip}\paragraph{選択肢の検討（査読待ち初稿）}\mbox{}\par

\noindent\textbf{1：不適合} 正誤正正。アは長いほど臨界荷重が小さくなるので誤。イを誤、ウ・エを正とする点も違う。\par

\noindent\textbf{2：不適合} 正正誤正。アは逆二乗則に反し、エは限界状態を一つだけに限定している。\par

\noindent\textbf{3：不適合} 誤誤誤正。イは弾性オイラー式の条件内では同じE・幾何・拘束で等しい座屈荷重となるため正。エも誤。\par

\noindent\textbf{4：不適合} 誤正正誤。ウの強度に安全率を掛けた許容応力は逆である。\par

\noindent\textbf{5：適合（前提内）} 誤正誤誤。アは長さ依存の方向が誤、イは弾性オイラー範囲で正、ウは除算でなく乗算のため誤、エは限界状態の単一化のため誤。\par

\paragraph{再利用する考え方} 材料の強さ、構造の安定、変形による機能低下を分ける。寸法だけを変えるとき何を固定したかを明示し、同じ理論でも候補番号を年度間で流用しない。

\paragraph{共通理論} \hyperref[theory:structural:buckling]{第\ref*{theory:structural:buckling}節}, \hyperref[theory:structural:allowable]{第\ref*{theory:structural:allowable}節}, \hyperref[theory:structural:limit]{第\ref*{theory:structural:limit}節}

\par 参考：\cite{struct_mit_beam,struct_mit_buckling,struct_mlit_limit}。

\paragraph{関連問題} \hyperref[q:h25-reg-basic-1-2]{H25-REG-BASIC-1-2}, \hyperref[q:r06-reg-basic-1-2]{R06-REG-BASIC-1-2}

\par 難易度 2（校準初稿）、優先度未評価。

\subsection{令和5年度 問1-2}
\label{q:r05-reg-basic-1-2}

\noindent\textbf{原典：} sources/R05/basic.pdf、PDF物理2頁。

\textbf{公式正答：} 4（single）。

\paragraph{問題の要旨} 細長い棒の軸圧縮による横変形を、静的座屈として説明する。

\Needspace{4\baselineskip}\paragraph{確認した条件・記述（独自要約）}\mbox{}\par

\noindent\textbf{判断方向：} 空欄の適切な語・数値又は正誤の組合せ。全候補の順序は各肢に記録。\par

\noindent\textbf{原典照合：} 16原典画像で17問の全候補、図、数値、不等号、表の列順を確認。\par

\noindent\textbf{前提：} 解法に記録した弾性範囲、荷重条件、幾何、独立性を明示する。\par

\paragraph{問われる核心} 細長い棒の軸圧縮による横変形を、静的座屈として説明する。

\paragraph{必要知識} 圧縮荷重、曲げ剛性EI、端部拘束に依存する有効長KL、座屈と共振の差。

\paragraph{考え方・正答を選ぶ理由（査読待ち初稿）} 軸圧縮のもとで横たわみを増やす効果が曲げの復元に勝つと直線状態が不安定になる。\(P_{\mathrm{cr}}=\pi^2EI/(KL)^2\) であり、断面円の \(I=\pi d^4/64\) より、小径ほど臨界荷重が小さい。図の現象はねじれや振動の共振ではなく曲げ座屈である。理想棒では臨界を超える条件を選び、実物の初期不整に対しては臨界前からたわみが生じ得る点も分ける。

\Needspace{4\baselineskip}\paragraph{選択肢の検討（査読待ち初稿）}\mbox{}\par

\noindent\textbf{1：不適合} 大きい・以下・ねじれ・共振。同じ長さ・材質では小径ほど曲げ剛性EIが小さく座屈しやすい。理想的な直線平衡の安定性を失うのは臨界荷重を超える側である。描かれた横変位は曲げであり、断面の軸回り回転であるねじれではない。静的圧縮での安定性喪失は座屈であり、周期的外力と固有振動の共振ではない。\par

\noindent\textbf{2：不適合} 小さい・超える・ねじれ・座屈。描かれた横変位は曲げであり、断面の軸回り回転であるねじれではない。\par

\noindent\textbf{3：不適合} 大きい・超える・曲げ・共振。同じ長さ・材質では小径ほど曲げ剛性EIが小さく座屈しやすい。静的圧縮での安定性喪失は座屈であり、周期的外力と固有振動の共振ではない。\par

\noindent\textbf{4：適合（前提内）} 小さい・超える・曲げ・座屈。小さい径・臨界荷重を超える・曲げ・座屈の対応である。\par

\noindent\textbf{5：不適合} 小さい・以下・曲げ・共振。理想的な直線平衡の安定性を失うのは臨界荷重を超える側である。静的圧縮での安定性喪失は座屈であり、周期的外力と固有振動の共振ではない。\par

\paragraph{再利用する考え方} 材料の強さ、構造の安定、変形による機能低下を分ける。寸法だけを変えるとき何を固定したかを明示し、同じ理論でも候補番号を年度間で流用しない。

\paragraph{共通理論} \hyperref[theory:structural:buckling]{第\ref*{theory:structural:buckling}節}

\par 参考：\cite{struct_mit_beam,struct_mit_buckling,struct_mlit_limit}。

\paragraph{関連問題} \hyperref[q:r01-reg-basic-1-4]{R01-REG-BASIC-1-4}

\par 難易度 2（校準初稿）、優先度未評価。

\section{安全率と費用モデルの最適化}

\subsection{令和元年度再試験 問1-5}
\label{q:r01-ret-basic-1-5}

\noindent\textbf{原典：} sources/R01\_retest/basic.pdf、PDF物理5頁。

\textbf{公式正答：} 2（single）。

\paragraph{問題の要旨} 安全率に依存する損害費用と対策費用の和を、定義域で最小化する。

\Needspace{4\baselineskip}\paragraph{確認した条件・記述（独自要約）}\mbox{}\par

\noindent\textbf{判断方向：} 空欄の適切な語・数値又は正誤の組合せ。全候補の順序は各肢に記録。\par

\noindent\textbf{原典照合：} 16原典画像で17問の全候補、図、数値、不等号、表の列順を確認。\par

\noindent\textbf{前提：} 解法に記録した弾性範囲、荷重条件、幾何、独立性を明示する。\par

\paragraph{問われる核心} 安全率に依存する損害費用と対策費用の和を、定義域で最小化する。

\paragraph{必要知識} 与えられた費用関数、微分による停留点、二階微分の正値、定義域と単位。

\paragraph{考え方・正答を選ぶ理由（査読待ち初稿）} 費用の単位を億円にそろえる。\(x\ge1\)、\(C(x)=90/(1+4x)+10x\) である。\(C^\prime=-360/(1+4x)^2+10=0\) より \(1+4x=6\)、\(x=5/4=1.25\) で候補2となる。\(C^{\prime\prime}=2880/(1+4x)^3>0\) なので許容域内の唯一の最小である。候補の費用は順に28、27.5、約27.857、28.75、30億円である。

\Needspace{4\baselineskip}\paragraph{選択肢の検討（査読待ち初稿）}\mbox{}\par

\noindent\textbf{1：不適合} 1.00。与式に代入した総費用は約28億円であり、最小の候補より大きい。\par

\noindent\textbf{2：適合（前提内）} 1.25。与式に代入した総費用は約27.5億円であり、連続域での停留点かつ最小である。\par

\noindent\textbf{3：不適合} 1.50。与式に代入した総費用は約27.857142857億円であり、最小の候補より大きい。\par

\noindent\textbf{4：不適合} 1.75。与式に代入した総費用は約28.75億円であり、最小の候補より大きい。\par

\noindent\textbf{5：不適合} 2.00。与式に代入した総費用は約30億円であり、最小の候補より大きい。\par

\paragraph{再利用する考え方} 材料の強さ、構造の安定、変形による機能低下を分ける。寸法だけを変えるとき何を固定したかを明示し、同じ理論でも候補番号を年度間で流用しない。

\paragraph{共通理論} \hyperref[theory:structural:optimization]{第\ref*{theory:structural:optimization}節}

\paragraph{関連問題} \hyperref[q:r04-reg-basic-1-4]{R04-REG-BASIC-1-4}

\par 難易度 2（校準初稿）、優先度未評価。

\subsection{令和4年度 問1-4}
\label{q:r04-reg-basic-1-4}

\noindent\textbf{原典：} sources/R04/basic.pdf、PDF物理3頁。

\textbf{公式正答：} 1（single）。

\paragraph{問題の要旨} 安全率に依存する損害費用と対策費用の和を、定義域で最小化する。

\Needspace{4\baselineskip}\paragraph{確認した条件・記述（独自要約）}\mbox{}\par

\noindent\textbf{判断方向：} 空欄の適切な語・数値又は正誤の組合せ。全候補の順序は各肢に記録。\par

\noindent\textbf{原典照合：} 16原典画像で17問の全候補、図、数値、不等号、表の列順を確認。\par

\noindent\textbf{前提：} 解法に記録した弾性範囲、荷重条件、幾何、独立性を明示する。\par

\paragraph{問われる核心} 安全率に依存する損害費用と対策費用の和を、定義域で最小化する。

\paragraph{必要知識} 与えられた費用関数、微分による停留点、二階微分の正値、定義域と単位。

\paragraph{考え方・正答を選ぶ理由（査読待ち初稿）} 費用の単位を億円にそろえる。\(x>1\)、\(C(x)=9/(1+x)+x\) である。\(C^\prime=-9/(1+x)^2+1=0\) より \(1+x=3\)、\(x=2\) で候補1となる。\(C^{\prime\prime}=18/(1+x)^3>0\) なので許容域で唯一の最小である。候補の費用は順に5、約5.071、5.25、5.5、5.8億円である。

\Needspace{4\baselineskip}\paragraph{選択肢の検討（査読待ち初稿）}\mbox{}\par

\noindent\textbf{1：適合（前提内）} 2.0。与式に代入した総費用は約5億円であり、連続域での停留点かつ最小である。\par

\noindent\textbf{2：不適合} 2.5。与式に代入した総費用は約5.071428571億円であり、最小の候補より大きい。\par

\noindent\textbf{3：不適合} 3.0。与式に代入した総費用は約5.25億円であり、最小の候補より大きい。\par

\noindent\textbf{4：不適合} 3.5。与式に代入した総費用は約5.5億円であり、最小の候補より大きい。\par

\noindent\textbf{5：不適合} 4.0。与式に代入した総費用は約5.8億円であり、最小の候補より大きい。\par

\paragraph{再利用する考え方} 材料の強さ、構造の安定、変形による機能低下を分ける。寸法だけを変えるとき何を固定したかを明示し、同じ理論でも候補番号を年度間で流用しない。

\paragraph{共通理論} \hyperref[theory:structural:optimization]{第\ref*{theory:structural:optimization}節}

\paragraph{関連問題} \hyperref[q:r01-ret-basic-1-5]{R01-RET-BASIC-1-5}

\par 難易度 2（校準初稿）、優先度未評価。

\section{崩壊・共振・圧壊と構造モデル}

\subsection{令和2年度 問1-3}
\label{q:r02-reg-basic-1-3}

\noindent\textbf{原典：} sources/R02/basic.pdf、PDF物理3頁。

\textbf{公式正答：} 3（single）。

\par\textbf{検討未確定：} 短柱の圧縮が破断強度を超えるというエの語は材料・圧縮強度・降伏の条件が曖昧。延性材料が一定の引張破断強度で圧縮破断するとは一般化しない。初回3/公式3を保持。

\paragraph{問題の要旨} 静的崩壊、共振、座屈、短柱の圧縮破損と構造解析のモデルを区別する。

\Needspace{4\baselineskip}\paragraph{確認した条件・記述（独自要約）}\mbox{}\par

\noindent\textbf{判断方向：} 空欄の適切な語・数値又は正誤の組合せ。全候補の順序は各肢に記録。\par

\noindent\textbf{原典照合：} 16原典画像で17問の全候補、図、数値、不等号、表の列順を確認。\par

\noindent\textbf{前提：} 解法に記録した弾性範囲、荷重条件、幾何、独立性を明示する。\par

\paragraph{問われる核心} 静的崩壊、共振、座屈、短柱の圧縮破損と構造解析のモデルを区別する。

\paragraph{必要知識} 塑性崩壊、外力周波数と固有振動数、細長比、圧縮材料破損、構造解析モデル。

\paragraph{考え方・正答を選ぶ理由（査読待ち初稿）} アの崩壊は塑性ヒンジ等による荷重支持機構の喪失であり弾性範囲だけで決めない。イの共振は外力の周波数が固有振動数へ近づくときに起き、応力が振動数へ達するという比較は次元も違う。ウの細長い柱は圧縮座屈する。エは短く太い柱で材料圧縮破損が支配し得るという趣旨、オは荷重・支持条件と部材を表す解析モデルである。従って誤誤正正正、3とする。ただしエの破断強度を全材料の圧縮破壊に同じように適用しない。

\Needspace{4\baselineskip}\paragraph{選択肢の検討（査読待ち初稿）}\mbox{}\par

\noindent\textbf{1：不適合} 正正正誤誤。アは塑性崩壊を弾性で決める点が誤り、イも応力と振動数を比べている。エ・オを誤りとする点も違う。\par

\noindent\textbf{2：不適合} 誤正正正誤。イの応力が固有振動数へ達する説明は単位も機構も違う。オの解析モデルの定義を否定できない。\par

\noindent\textbf{3：適合（前提内）} 誤誤正正正。アは誤、イは誤、ウは圧縮座屈で正、エは短柱の材料圧縮破損を意図した正、オは部材と力・拘束のモデルで正となる。\par

\noindent\textbf{4：不適合} 正誤誤正正。アは塑性機構の崩壊荷重であって弾性だけでは決まらず、ウの圧縮座屈は正である。\par

\noindent\textbf{5：不適合} 正正誤誤正。ア・イは誤、ウは正なので最初の三記号が一致しない。\par

\paragraph{再利用する考え方} 材料の強さ、構造の安定、変形による機能低下を分ける。寸法だけを変えるとき何を固定したかを明示し、同じ理論でも候補番号を年度間で流用しない。

\paragraph{共通理論} \hyperref[theory:structural:modes]{第\ref*{theory:structural:modes}節}, \hyperref[theory:structural:buckling]{第\ref*{theory:structural:buckling}節}

\par 参考：\cite{struct_mit_beam,struct_mit_buckling,struct_mlit_limit}。

\paragraph{関連問題} \hyperref[q:h24-reg-basic-1-2]{H24-REG-BASIC-1-2}

\par 難易度 2（校準初稿）、優先度未評価。

\section{疲労・クリープと局所応力集中}

\subsection{令和4年度 問1-1}
\label{q:r04-reg-basic-1-1}

\noindent\textbf{原典：} sources/R04/basic.pdf、PDF物理1頁。

\textbf{公式正答：} 1（single）。

\par\textbf{検討未確定：} 疲労の平均応力依存は振幅等の同条件比較を要する。急変部の局所von Mises値は荷重/形状/拘束に依存し、常に上がる全称の証明ではない。公式/初回1を維持し条件を限定。

\paragraph{問題の要旨} 疲労の平均応力、クリープ温度、曲げ剛性と応力集中を区別する。

\Needspace{4\baselineskip}\paragraph{確認した条件・記述（独自要約）}\mbox{}\par

\noindent\textbf{判断方向：} 空欄の適切な語・数値又は正誤の組合せ。全候補の順序は各肢に記録。\par

\noindent\textbf{原典照合：} 16原典画像で17問の全候補、図、数値、不等号、表の列順を確認。\par

\noindent\textbf{前提：} 解法に記録した弾性範囲、荷重条件、幾何、独立性を明示する。\par

\paragraph{問われる核心} 疲労の平均応力、クリープ温度、曲げ剛性と応力集中を区別する。

\paragraph{必要知識} 応力振幅/平均応力、通常の平均応力補正の条件、高温クリープ、EIと変位、局所応力と公称応力。

\paragraph{考え方・正答を選ぶ理由（査読待ち初稿）} Aは通常の疲労限度線図の引張平均応力側である。同じ応力振幅なら平均引張応力の増加は一般に疲労寿命へ不利となるが、材料・残留応力・塑性域へ無条件には拡張しない。Bは高温での時間依存変形、Cは同じ荷重と断面の \(\delta\propto E^{-1}\) より小さいヤング率、Dは断面急変部の局所応力の上昇である。従って1を選ぶ。von Mises応力の局所値は荷重と形状に依存し、急変させればどの点でも必ず高くなるという命題にはしない。

\Needspace{4\baselineskip}\paragraph{選択肢の検討（査読待ち初稿）}\mbox{}\par

\noindent\textbf{1：適合（前提内）} 引張・高い・小さい・高い。振幅等を固定した通常の金属疲労モデルでは引張平均応力が不利、高温でクリープ、小Eで大変位、急変部で局所応力が高い。\par

\noindent\textbf{2：不適合} 引張・高い・大きい・小さい。同じ荷重と幾何では変位はEの逆数であり、小Eが大変位である。応力集中は局所上昇を表す。\par

\noindent\textbf{3：不適合} 圧縮・高い・小さい・小さい。一般的平均応力モデルの不利側は引張であり、応力集中を局所低下とはしない。\par

\noindent\textbf{4：不適合} 圧縮・高強度・小さい・高い。Aは引張側、Bは高温である。高強度を温度条件の代わりにはできない。\par

\noindent\textbf{5：不適合} 引張・高強度・大きい・高い。クリープのBは高温であり、曲げ変位のCは小さいEである。\par

\paragraph{再利用する考え方} 材料の強さ、構造の安定、変形による機能低下を分ける。寸法だけを変えるとき何を固定したかを明示し、同じ理論でも候補番号を年度間で流用しない。

\paragraph{共通理論} \hyperref[theory:structural:fatigue]{第\ref*{theory:structural:fatigue}節}, \hyperref[theory:structural:creep]{第\ref*{theory:structural:creep}節}, \hyperref[theory:structural:beam]{第\ref*{theory:structural:beam}節}, \hyperref[theory:structural:concentration]{第\ref*{theory:structural:concentration}節}

\par 参考：\cite{struct_nims_data,struct_nasa_mean,struct_mit_beam}。

\par 難易度 2（校準初稿）、優先度未評価。

\section{使用・終局・疲労の限界状態}

\subsection{平成24年度 問1-2}
\label{q:h24-reg-basic-1-2}

\noindent\textbf{原典：} sources/H24/basic.pdf、PDF物理2頁。

\textbf{公式正答：} 5（single）。

\par\textbf{検討未確定：} 全ての構成部材の破壊まで至らなければ安全性が十分という前文は、機構形成や全体座屈を除外しておらず一般には保証できない。公式5と初回5を保持し、余力の条件を限定する。

\paragraph{問題の要旨} 不静定構造の余力、部分損傷、限界状態と非線形解析の関係を整理する。

\Needspace{4\baselineskip}\paragraph{確認した条件・記述（独自要約）}\mbox{}\par

\noindent\textbf{判断方向：} 空欄の適切な語・数値又は正誤の組合せ。全候補の順序は各肢に記録。\par

\noindent\textbf{原典照合：} 16原典画像で17問の全候補、図、数値、不等号、表の列順を確認。\par

\noindent\textbf{前提：} 解法に記録した弾性範囲、荷重条件、幾何、独立性を明示する。\par

\paragraph{問われる核心} 不静定構造の余力、部分損傷、限界状態と非線形解析の関係を整理する。

\paragraph{必要知識} 静定と不静定、荷重再配分、塑性ヒンジ、材料・幾何非線形、終局/使用/疲労の限界。

\paragraph{考え方・正答を選ぶ理由（査読待ち初稿）} Aを高い、Bを限界状態、Cを部分的、Dを非線形性とする5である。余剰拘束のある構造で一部が降伏しても荷重を他部へ再配分できる場合がある。しかし塑性機構、全体座屈、接合部破壊が先に生じ得るため、「全ての部材が破壊しなければ安全」という前文を普遍命題にしない。

\Needspace{4\baselineskip}\paragraph{選択肢の検討（査読待ち初稿）}\mbox{}\par

\noindent\textbf{1：不適合} 高い・限界状態・部分的・線形性。損傷後の応答を線形のままとするDが違う。\par

\noindent\textbf{2：不適合} 低い・限界状態・全体的・非線形性。Aは余剰拘束が多い高不静定、Cは部分損傷を許す説明なので全体的とする点が違う。\par

\noindent\textbf{3：不適合} 高い・弾性・部分的・線形性。弾性設計に限定すると部分損傷後の余力を捉えず、Dも非線形である。\par

\noindent\textbf{4：不適合} 低い・弾性・全体的・線形性。高不静定の余力を低不静定へ置換し、弾性・全体的・線形へ限定している。\par

\noindent\textbf{5：適合（前提内）} 高い・限界状態・部分的・非線形性。想定された延性と荷重再配分のモデルでは高不静定、限界状態、部分的損傷、非線形が対応する。ただし全構成部材が破壊しないだけでは全体安全の証明にならない。\par

\paragraph{再利用する考え方} 材料の強さ、構造の安定、変形による機能低下を分ける。寸法だけを変えるとき何を固定したかを明示し、同じ理論でも候補番号を年度間で流用しない。

\paragraph{共通理論} \hyperref[theory:structural:limit]{第\ref*{theory:structural:limit}節}, \hyperref[theory:structural:redundancy]{第\ref*{theory:structural:redundancy}節}

\par 参考：\cite{struct_mit_beam,struct_mit_buckling,struct_mlit_limit}。

\paragraph{関連問題} \hyperref[q:r02-reg-basic-1-3]{R02-REG-BASIC-1-3}

\par 難易度 2（校準初稿）、優先度未評価。

\subsection{令和6年度 問1-2}
\label{q:r06-reg-basic-1-2}

\noindent\textbf{原典：} sources/R06/basic.pdf、PDF物理2頁。

\textbf{公式正答：} 2（single）。

\paragraph{問題の要旨} 通常使用、繰返し荷重、最大耐荷性能の三つの要求を限界状態へ対応させる。

\Needspace{4\baselineskip}\paragraph{確認した条件・記述（独自要約）}\mbox{}\par

\noindent\textbf{判断方向：} 空欄の適切な語・数値又は正誤の組合せ。全候補の順序は各肢に記録。\par

\noindent\textbf{原典照合：} 16原典画像で17問の全候補、図、数値、不等号、表の列順を確認。\par

\noindent\textbf{前提：} 解法に記録した弾性範囲、荷重条件、幾何、独立性を明示する。\par

\paragraph{問われる核心} 通常使用、繰返し荷重、最大耐荷性能の三つの要求を限界状態へ対応させる。

\paragraph{必要知識} 使用性、疲労による累積損傷、終局耐荷性能、規準別の分類粒度。

\paragraph{考え方・正答を選ぶ理由（査読待ち初稿）} 過大なたわみや振動等で通常の使用に支障が生じる状態は使用限界である。繰返し荷重の損傷が疲労限界、崩壊や不安定等で最大耐荷性能を失う状態が終局限界である。使用・疲労・終局の2を選ぶ。基準によっては疲労を終局に含める分類もあるが、本問の三項目対応を変えない。

\Needspace{4\baselineskip}\paragraph{選択肢の検討（査読待ち初稿）}\mbox{}\par

\noindent\textbf{1：不適合} 使用・終局・疲労。Bの繰返し荷重は疲労、Cの最大耐荷は終局であり入れ替わっている。\par

\noindent\textbf{2：適合（前提内）} 使用・疲労・終局。Aは通常の機能・快適性等の使用限界、Bは疲労、Cは終局であり対応する。\par

\noindent\textbf{3：不適合} 終局・疲労・使用。Aは使用、Cは終局なので逆転している。\par

\noindent\textbf{4：不適合} 疲労・使用・終局。Aは使用、Bは繰返しの疲労なので逆転している。\par

\noindent\textbf{5：不適合} 疲労・終局・使用。三項目とも用途との対応が一致しない。\par

\paragraph{再利用する考え方} 材料の強さ、構造の安定、変形による機能低下を分ける。寸法だけを変えるとき何を固定したかを明示し、同じ理論でも候補番号を年度間で流用しない。

\paragraph{共通理論} \hyperref[theory:structural:limit]{第\ref*{theory:structural:limit}節}

\par 参考：\cite{struct_mlit_limit}。

\paragraph{関連問題} \hyperref[q:r03-reg-basic-1-5]{R03-REG-BASIC-1-5}

\par 難易度 2（校準初稿）、優先度未評価。

\section{強度・荷重の同時分布と破壊確率}

\subsection{平成29年度 問1-6}
\label{q:h29-reg-basic-1-6}

\noindent\textbf{原典：} sources/H29/basic.pdf、PDF物理5頁。

\textbf{公式正答：} 2（single）。

\par\textbf{検討未確定：} 強度と荷重の独立性が原文に明記されない。周辺一様分布だけでは破壊確率を一意に決定できない。独立なら初回2/公式2、同じ周辺を持つ別結合で0又は2/15となる反例を保存する。

\paragraph{問題の要旨} 強度と荷重が一様分布するモデルで、荷重が強度を超える確率と独立性の役割を検討する。

\Needspace{4\baselineskip}\paragraph{確認した条件・記述（独自要約）}\mbox{}\par

\noindent\textbf{判断方向：} 空欄の適切な語・数値又は正誤の組合せ。全候補の順序は各肢に記録。\par

\noindent\textbf{原典照合：} 16原典画像で17問の全候補、図、数値、不等号、表の列順を確認。\par

\noindent\textbf{前提：} 解法に記録した弾性範囲、荷重条件、幾何、独立性を明示する。\par

\paragraph{問われる核心} 強度と荷重が一様分布するモデルで、荷重が強度を超える確率と独立性の役割を検討する。

\paragraph{必要知識} 同時分布と周辺分布、独立性、破壊事象S>R、密度の面積積分。

\paragraph{考え方・正答を選ぶ理由（査読待ち初稿）} 強度を \(R\sim U[18,23]\)、荷重を \(S\sim U[10,20]\) とする。独立なら \(f_{R,S}=1/(5\cdot10)\)。破壊条件 \(S>R\) の領域は \(18\le R\le20,\ R<S\le20\) であるから、\(P(S>R)=\int_{18}^{20}(20-r)/50\,dr=4/100=0.04\)、候補2である。ただし原文には独立の明記がない。共通の \(U\sim U[0,1]\) を用い、\(R=18+5U,S=10+10U\) なら破壊確率は0、\(R=18+5U,S=20-10U\) なら \(U<2/15\) より確率2/15である。いずれも同じ周辺一様分布を持つ。従って公式0.04は独立モデルを補った条件付きの解として示す。

\Needspace{4\baselineskip}\paragraph{選択肢の検討（査読待ち初稿）}\mbox{}\par

\noindent\textbf{1：不適合} 0.02。独立モデルの三角形面積をさらに半分にすると0.02となるが、その二重の半減は不要である。\par

\noindent\textbf{2：適合（前提内）} 0.04。独立を仮定すれば支持矩形の面積50に対し破壊三角形面積2なので0.04となる。\par

\noindent\textbf{3：不適合} 0.08。一辺2の正方形を全て破壊領域にすると0.08となる。対角線の半分だけが該当する。\par

\noindent\textbf{4：不適合} 0.10。0.10は独立モデルの領域積分と一致しない。周辺分布だけで確率が決まる値でもない。\par

\noindent\textbf{5：不適合} 0.20。強度の区間長5の逆数0.20は密度であり、破壊事象の確率ではない。\par

\paragraph{再利用する考え方} 材料の強さ、構造の安定、変形による機能低下を分ける。寸法だけを変えるとき何を固定したかを明示し、同じ理論でも候補番号を年度間で流用しない。

\paragraph{共通理論} \hyperref[theory:structural:reliability]{第\ref*{theory:structural:reliability}節}

\par 難易度 4（校準初稿）、優先度未評価。

\section{比強度と材料価格当たり性能}

\subsection{令和5年度 問1-1}
\label{q:r05-reg-basic-1-1}

\noindent\textbf{原典：} sources/R05/basic.pdf、PDF物理1頁。

\textbf{公式正答：} 5（single）。

\paragraph{問題の要旨} 与えられた鋼とCFRPの密度・強度・価格から、比強度とその価格当たり性能を比較する。

\Needspace{4\baselineskip}\paragraph{確認した条件・記述（独自要約）}\mbox{}\par

\noindent\textbf{判断方向：} 空欄の適切な語・数値又は正誤の組合せ。全候補の順序は各肢に記録。\par

\noindent\textbf{原典照合：} 16原典画像で17問の全候補、図、数値、不等号、表の列順を確認。\par

\noindent\textbf{前提：} 解法に記録した弾性範囲、荷重条件、幾何、独立性を明示する。\par

\paragraph{問われる核心} 与えられた鋼とCFRPの密度・強度・価格から、比強度とその価格当たり性能を比較する。

\paragraph{必要知識} 比強度σ/ρ、同じ質量の荷重支持能力、価格/比強度、所与価格と方向条件。

\paragraph{考え方・正答を選ぶ理由（査読待ち初稿）} 比強度は \(\sigma/\rho\) である。与えられた数値だけを使うと、CFRP/鋼の比強度比は \((2000/1600)/(400/7900)=24.6875\)、kg当たり価格比は \(16000/60=266.6667\)。単位比強度当たり価格比は \((16000/(2000/1600))/(60/(400/7900))\simeq10.8017\) で、最も近い約10、5を選ぶ。比強度が大きいことと費用当たり有利なことを区別する。価格は設問所与であり現在の相場ではない。CFRPの方向依存や圧縮特性を所与の引張強度へ無断で補わない。

\Needspace{4\baselineskip}\paragraph{選択肢の検討（査読待ち初稿）}\mbox{}\par

\noindent\textbf{1：不適合} 強度/密度・CFRP・2。比強度の定義と優れる材料は合うが、単位比強度当たり価格の倍率は約10.8であり2ではない。\par

\noindent\textbf{2：不適合} 密度/強度・CFRP・10。密度/強度は比強度の逆数である。\par

\noindent\textbf{3：不適合} 密度/強度・鋼・2。定義を逆転し、与条件での優れる材料と価格倍率も違う。\par

\noindent\textbf{4：不適合} 強度/密度・鋼・2。定義は合うが、強度/密度はCFRPの方が大きく、価格倍率も2ではない。\par

\noindent\textbf{5：適合（前提内）} 強度/密度・CFRP・10。強度/密度でCFRPが優れ、価格当たりの比較ではCFRPの単位比強度当たり費用が約10倍となる。\par

\paragraph{再利用する考え方} 材料の強さ、構造の安定、変形による機能低下を分ける。寸法だけを変えるとき何を固定したかを明示し、同じ理論でも候補番号を年度間で流用しない。

\paragraph{共通理論} \hyperref[theory:structural:specific]{第\ref*{theory:structural:specific}節}

\par 難易度 3（校準初稿）、優先度未評価。

\section{引張試験と公称応力・弾塑性}

\subsection{平成29年度 問1-4}
\label{q:h29-reg-basic-1-4}

\noindent\textbf{原典：} sources/H29/basic.pdf、PDF物理3頁。

\textbf{公式正答：} 5（single）。

\par\textbf{検討未確定：} 前文は荷重を断面積で割った値を応力とするが初期断面積を明示しない。最大値を引張強さとするには公称応力を用いる条件が必要。真応力の最大と同一視せず初回/公式5を保持。

\paragraph{問題の要旨} 引張試験の測定値と、弾性・塑性・引張強さの定義を対応させる。

\Needspace{4\baselineskip}\paragraph{確認した条件・記述（独自要約）}\mbox{}\par

\noindent\textbf{判断方向：} 空欄の適切な語・数値又は正誤の組合せ。全候補の順序は各肢に記録。\par

\noindent\textbf{原典照合：} 16原典画像で17問の全候補、図、数値、不等号、表の列順を確認。\par

\noindent\textbf{前提：} 解法に記録した弾性範囲、荷重条件、幾何、独立性を明示する。\par

\paragraph{問われる核心} 引張試験の測定値と、弾性・塑性・引張強さの定義を対応させる。

\paragraph{必要知識} 初期断面積と初期標点距離、荷重/伸び、弾性/塑性、延性と公称引張強さ。

\paragraph{考え方・正答を選ぶ理由（査読待ち初稿）} 測定する荷重 \(F\) と伸び \(\Delta L\) から、公称応力 \(F/A_0\) と公称ひずみ \(\Delta L/L_0\) を求める。除荷で回復する変形が弾性、残る変形が塑性である。単調引張試験での \(\max F/A_0\) が引張強さであり、破断時の応力、繰返し荷重での疲労限度、真応力 \(F/A(t)\) と分ける。従って伸び・弾性・塑性・引張強さを選ぶ。

\Needspace{4\baselineskip}\paragraph{選択肢の検討（査読待ち初稿）}\mbox{}\par

\noindent\textbf{1：不適合} ひずみ・弾性・延性・疲労限度。直接測る長さの増分は伸びであり、ひずみは初期長さで除した値である。残る変形の性質は塑性であり、延性は破断まで塑性変形を許す能力である。単調引張の最大公称応力は引張強さであり、疲労限度や破断時応力と区別する。\par

\noindent\textbf{2：不適合} 伸び・塑性・弾性・引張強さ。除荷後に回復する性質は弾性である。残る変形の性質は塑性であり、延性は破断まで塑性変形を許す能力である。\par

\noindent\textbf{3：不適合} 伸び・弾性・延性・疲労限度。残る変形の性質は塑性であり、延性は破断まで塑性変形を許す能力である。単調引張の最大公称応力は引張強さであり、疲労限度や破断時応力と区別する。\par

\noindent\textbf{4：不適合} ひずみ・延性・塑性・破断強度。直接測る長さの増分は伸びであり、ひずみは初期長さで除した値である。除荷後に回復する性質は弾性である。単調引張の最大公称応力は引張強さであり、疲労限度や破断時応力と区別する。\par

\noindent\textbf{5：適合（前提内）} 伸び・弾性・塑性・引張強さ。伸び・弾性・塑性・引張強さが対応する。公称応力の断面積は初期値である。\par

\paragraph{再利用する考え方} 材料の強さ、構造の安定、変形による機能低下を分ける。寸法だけを変えるとき何を固定したかを明示し、同じ理論でも候補番号を年度間で流用しない。

\paragraph{共通理論} \hyperref[theory:structural:tensile]{第\ref*{theory:structural:tensile}節}

\par 参考：\cite{struct_mit_beam,struct_mit_buckling,struct_mlit_limit}。

\paragraph{関連問題} \hyperref[q:r05-reg-basic-1-3]{R05-REG-BASIC-1-3}

\par 難易度 2（校準初稿）、優先度未評価。

\subsection{令和5年度 問1-3}
\label{q:r05-reg-basic-1-3}

\noindent\textbf{原典：} sources/R05/basic.pdf、PDF物理3頁。

\textbf{公式正答：} 3（single）。

\paragraph{問題の要旨} 引張試験の測定値と、弾性・塑性・引張強さの定義を対応させる。

\Needspace{4\baselineskip}\paragraph{確認した条件・記述（独自要約）}\mbox{}\par

\noindent\textbf{判断方向：} 空欄の適切な語・数値又は正誤の組合せ。全候補の順序は各肢に記録。\par

\noindent\textbf{原典照合：} 16原典画像で17問の全候補、図、数値、不等号、表の列順を確認。\par

\noindent\textbf{前提：} 解法に記録した弾性範囲、荷重条件、幾何、独立性を明示する。\par

\paragraph{問われる核心} 引張試験の測定値と、弾性・塑性・引張強さの定義を対応させる。

\paragraph{必要知識} 初期断面積と初期標点距離、荷重/伸び、弾性/塑性、延性と公称引張強さ。

\paragraph{考え方・正答を選ぶ理由（査読待ち初稿）} 測定する荷重 \(F\) と伸び \(\Delta L\) から、公称応力 \(F/A_0\) と公称ひずみ \(\Delta L/L_0\) を求める。除荷で回復する変形が弾性、残る変形が塑性である。単調引張試験での \(\max F/A_0\) が引張強さであり、破断時の応力、繰返し荷重での疲労限度、真応力 \(F/A(t)\) と分ける。従って伸び・弾性・塑性・引張強さを選ぶ。

\Needspace{4\baselineskip}\paragraph{選択肢の検討（査読待ち初稿）}\mbox{}\par

\noindent\textbf{1：不適合} ひずみ・弾性・延性・疲労限度。直接測る長さの増分は伸びであり、ひずみは初期長さで除した値である。残る変形の性質は塑性であり、延性は破断まで塑性変形を許す能力である。単調引張の最大公称応力は引張強さであり、疲労限度や破断時応力と区別する。\par

\noindent\textbf{2：不適合} 伸び・塑性・弾性・引張強さ。除荷後に回復する性質は弾性である。残る変形の性質は塑性であり、延性は破断まで塑性変形を許す能力である。\par

\noindent\textbf{3：適合（前提内）} 伸び・弾性・塑性・引張強さ。伸び・弾性・塑性・引張強さが対応する。公称応力の断面積は初期値である。\par

\noindent\textbf{4：不適合} 伸び・弾性・延性・疲労限度。残る変形の性質は塑性であり、延性は破断まで塑性変形を許す能力である。単調引張の最大公称応力は引張強さであり、疲労限度や破断時応力と区別する。\par

\noindent\textbf{5：不適合} ひずみ・延性・塑性・引張強さ。直接測る長さの増分は伸びであり、ひずみは初期長さで除した値である。除荷後に回復する性質は弾性である。\par

\paragraph{再利用する考え方} 材料の強さ、構造の安定、変形による機能低下を分ける。寸法だけを変えるとき何を固定したかを明示し、同じ理論でも候補番号を年度間で流用しない。

\paragraph{共通理論} \hyperref[theory:structural:tensile]{第\ref*{theory:structural:tensile}節}

\par 参考：\cite{struct_mit_beam,struct_mit_buckling,struct_mlit_limit}。

\paragraph{関連問題} \hyperref[q:h29-reg-basic-1-4]{H29-REG-BASIC-1-4}

\par 難易度 2（校準初稿）、優先度未評価。

\section{ユニバーサルデザイン}

\subsection{平成24年度 問1-4}
\label{q:h24-reg-basic-1-4}

\noindent\textbf{原典：} sources/H24/basic.pdf、PDF物理4頁。

\textbf{公式正答：} 5（single）。

\par\textit{原典本文の校正と独自解説は未作成。分類のみの登録項目。}

\subsection{平成26年度 問1-1}
\label{q:h26-reg-basic-1-1}

\noindent\textbf{原典：} sources/H26/basic.pdf、PDF物理1頁。

\textbf{公式正答：} 5（single）。

\par\textit{原典本文の校正と独自解説は未作成。分類のみの登録項目。}

\subsection{平成30年度 問1-3}
\label{q:h30-reg-basic-1-3}

\noindent\textbf{原典：} sources/H30/basic.pdf、PDF物理3頁。

\textbf{公式正答：} 2（single）。

\par\textit{原典本文の校正と独自解説は未作成。分類のみの登録項目。}

\subsection{令和2年度 問1-1}
\label{q:r02-reg-basic-1-1}

\noindent\textbf{原典：} sources/R02/basic.pdf、PDF物理1頁。

\textbf{公式正答：} 3（single）。

\par\textit{原典本文の校正と独自解説は未作成。分類のみの登録項目。}

\subsection{令和3年度 問1-1}
\label{q:r03-reg-basic-1-1}

\noindent\textbf{原典：} sources/R03/basic.pdf、PDF物理1頁。

\textbf{公式正答：} 4（single）。

\par\textit{原典本文の校正と独自解説は未作成。分類のみの登録項目。}

\subsection{令和6年度 問1-1}
\label{q:r06-reg-basic-1-1}

\noindent\textbf{原典：} sources/R06/basic.pdf、PDF物理1頁。

\textbf{公式正答：} 1（single）。

\par\textit{原典本文の校正と独自解説は未作成。分類のみの登録項目。}

\section{未確定}

\subsection{平成25年度 問1-1}
\label{q:h25-reg-basic-1-1}

\noindent\textbf{原典：} sources/H25/basic.pdf、PDF物理1頁。

\textbf{公式正答：} 3（single）。

\par\textit{原典本文の校正と独自解説は未作成。分類のみの登録項目。}

\subsection{平成27年度 問1-4}
\label{q:h27-reg-basic-1-4}

\noindent\textbf{原典：} sources/H27/basic.pdf、PDF物理3頁。

\textbf{公式正答：} 3（single）。

\par\textit{原典本文の校正と独自解説は未作成。分類のみの登録項目。}

\subsection{平成30年度 問1-6}
\label{q:h30-reg-basic-1-6}

\noindent\textbf{原典：} sources/H30/basic.pdf、PDF物理5頁。

\textbf{公式正答：} 4（single）。

\par\textit{原典本文の校正と独自解説は未作成。分類のみの登録項目。}

\subsection{令和元年度 問1-2}
\label{q:r01-reg-basic-1-2}

\noindent\textbf{原典：} sources/R01/basic.pdf、PDF物理2頁。

\textbf{公式正答：} 2（single）。

\par\textit{原典本文の校正と独自解説は未作成。分類のみの登録項目。}

\subsection{令和3年度 問1-6}
\label{q:r03-reg-basic-1-6}

\noindent\textbf{原典：} sources/R03/basic.pdf、PDF物理5頁。

\textbf{公式正答：} 3（single）。

\par\textit{原典本文の校正と独自解説は未作成。分類のみの登録項目。}

\subsection{令和4年度 問1-2}
\label{q:r04-reg-basic-1-2}

\noindent\textbf{原典：} sources/R04/basic.pdf、PDF物理2頁。

\textbf{公式正答：} 4（single）。

\par\textit{原典本文の校正と独自解説は未作成。分類のみの登録項目。}

\subsection{令和4年度 問1-6}
\label{q:r04-reg-basic-1-6}

\noindent\textbf{原典：} sources/R04/basic.pdf、PDF物理4頁。

\textbf{公式正答：} 2（single）。

\par\textit{原典本文の校正と独自解説は未作成。分類のみの登録項目。}

\subsection{令和6年度 問1-4}
\label{q:r06-reg-basic-1-4}

\noindent\textbf{原典：} sources/R06/basic.pdf、PDF物理3頁。

\textbf{公式正答：} 3（single）。

\par\textit{原典本文の校正と独自解説は未作成。分類のみの登録項目。}

\subsection{令和6年度 問1-6}
\label{q:r06-reg-basic-1-6}

\noindent\textbf{原典：} sources/R06/basic.pdf、PDF物理5頁。

\textbf{公式正答：} 3（single）。

\par\textit{原典本文の校正と独自解説は未作成。分類のみの登録項目。}
